\chapter{Measurement Validity: Accuracy, Neutrality, and Noise}
\label{ch:validity}

A characterization inherits the validity of its measurements. This chapter
establishes, in order: that the system under measurement was computing
\emph{correct} trajectories (\S\ref{sec:val-accuracy}); that observing it did
not change its behaviour (\S\ref{sec:val-neutrality}); and that residual
run-to-run variation is understood and bounded
(\S\ref{sec:val-nondeterminism}). Together these answer RQ5 and constitute
finding \textbf{F13} of the evidence chain.

\section{Accuracy validation: the 141-configuration matrix}
\label{sec:val-accuracy}

Every combination of dataset, sensor modality, and pipeline mode that the
on-disk corpora support was materialized into a configuration---141 in
total---and each was run to completion and scored against ground truth with
the vendor's own metrics (RMSE APE after alignment; segment-relative RTE;
rotational RE). A convergence gate (segment-relative translational drift
below 5\%, a scale-independent criterion valid from room-scale to
kilometre-scale trajectories) separates runs that track from runs that
diverge; every excluded run is itemized with a stated cause rather than
silently dropped.

Table~\ref{tab:accuracy} compares the converged mode averages against the
cuVSLAM technical report~\cite{cuvslam2025}. Two facts carry the chapter.
First, \emph{the modes used for the entire memory characterization---stereo
and RGB-D SLAM---reproduce the published accuracy}: EuRoC stereo to within
millimetres, TUM RGB-D better than published, and KITTI matching on the
definition-stable per-segment drift metric (0.82\% at the 500\,m segment
versus the 0.85\% leaderboard figure; absolute APE over kilometre paths is
scale-sensitive and reported as directional only). Second, every
discrepancy in the remaining modes has a named, verified cause, none of
which touches the characterized configurations.

\begin{table}[t]
\centering
\caption{Converged mode-average accuracy vs.\ the cuVSLAM technical report
(RMSE APE in metres, odometry/SLAM). The characterization uses the stereo
and RGB-D modes.}
\label{tab:accuracy}
\begin{tabular}{@{}lccl@{}}
\toprule
mode & ours & published & verdict \\
\midrule
EuRoC stereo & 0.114 / 0.051 & 0.13 / 0.054 & millimetre-level match \\
TUM RGB-D & 0.060 / 0.047 & 0.11 / 0.065 & better than published \\
ICL-NUIM (mono-depth) & 0.099 / 0.136 & 0.026 & same order (cm) \\
KITTI stereo & 4.01 / 3.63 & 3.00 / 1.98 & segment drift matches
(0.82\% vs 0.85\%) \\
EuRoC stereo-inertial & 0.401 / 0.328 & 0.19 / 0.13 & under-tuned IMU
(disclosed) \\
\bottomrule
\end{tabular}
\end{table}

\subsection{Discrepancies and their causes}
\label{sec:val-discrepancies}

\textbf{Inertial mode is under-tuned---by configuration, not by cuVSLAM.}
Only a minority of EuRoC inertial runs converge, and---diagnostically---
enabling the IMU makes accuracy \emph{worse} than pure stereo on the same
sequences. That inversion is the classic signature of mis-scaled IMU noise
densities or camera--IMU extrinsics: the estimator over-trusts a mis-modelled
sensor. The vendor calibrates these per dataset; the matrix used a generic
inertial configuration. The defect is confined to a mode the memory
characterization never uses, and the remedy (populating per-sequence IMU
parameters from the datasets' calibration files) is configuration work,
tracked in the roadmap.

\textbf{The hardest EuRoC sequence diverges in every mode.} V2\_03 combines
aggressive motion with heavy motion blur; the vendor's own report publishes
it only under stereo-inertial and excludes it from stereo averages. The same
exclusion is applied here, with the sequence retained in per-sequence
tables.

\textbf{TUM-VI raw fisheye is invalid input, not a tracking failure.}
TUM-VI's $\sim$195$^\circ$ fisheye exceeds the pinhole-undistortion field of
view the library supports ($<$180$^\circ$); feeding raw fisheye produces
kilometre-scale divergence deterministically. The documented preparation
(undistort to pinhole with edge masks) was identified in the dataset notes;
runs without it are classified invalid-input and excluded.

\textbf{Monocular needs scale-free alignment.} Monocular SLAM recovers shape
but not metric scale; scoring it with SE(3) alignment conflates scale error
with drift. Monocular rows are therefore excluded from SE(3) comparisons
pending Sim(3) evaluation~\cite{umeyama1991}---an evaluator-metric choice,
not a tracking result.

\section{Profiling is accuracy-neutral}
\label{sec:val-neutrality}

The deepest reflexivity threat is that instrumentation changes the system it
observes. Three experiments close it.

\textbf{Instrumented build equals baseline build (QoR).} The TaggedAllocator
instrumentation (Chapter~\ref{ch:attribution}) lives in a rebuilt cuVSLAM
binary. Re-running representative configurations on both wheels shows the
instrumented build's trajectories equal the baseline's---bit-identical on
EuRoC configurations, and within run-to-run nondeterminism elsewhere
(largest observed difference 0.24\,m over a kilometre-scale KITTI path,
within the sequence's own re-run scatter; see
\S\ref{sec:val-nondeterminism}). The journaling itself is gated by an
environment variable and touches only allocation paths, which execute at
initialization, not per frame.

\textbf{Each profiler is trajectory-neutral.} Running the same configuration
plain and under each instrument shows: nsys and ncu produce
\emph{bit-identical} trajectories on deterministic (synchronous) modes, and
NVBit-traced runs agree within $\sim$2\,mm APE---the instrumentation slows
execution greatly but does not alter the computation.

\textbf{Neutrality holds across the full feature matrix.} A 192-variant
campaign (every accuracy configuration profiled as-is, plus finer feature
toggles---sync/async/CPU SLAM, motion-model and denoising toggles,
landmark export, multi-camera precision, unrectified input, depth-stereo
tracking---on representative sequences per modality) compares profiled
against plain APE with tolerance $\max(5\,\text{cm},\,5\%)$: \textbf{166/192
pass outright}, with 121 bit-identical. Every one of the 26 flagged variants
was investigated and traced to a benign cause---dominated by the
nondeterministic scatter of long loop-closing sequences
(\S\ref{sec:val-nondeterminism}) and by toggles that legitimately change the
computation (e.g.\ CPU-SLAM changes the back-end, hence the trajectory,
independent of profiling); a re-check against plain re-runs disproved the
one suspected genuine interference case.

\section{Nondeterminism, characterized}
\label{sec:val-nondeterminism}

GPU SLAM is not perfectly repeatable: floating-point reductions reorder
across runs, and asynchronous modes add scheduling freedom. Rather than
averaging this away, the campaign measured it. On short and mid-length
sequences, repeated runs are bit-identical in synchronous modes. On long
loop-closing sequences (kitti00 is the canonical case) the final trajectory
is \emph{bimodal}: loop-closure acceptance is threshold-sensitive, and a
borderline loop either fires or does not, moving whole-path APE by tenths
of a metre. The decisive control: \emph{plain, uninstrumented re-runs
scatter by the same amount and between the same modes}. The scatter is a
property of the workload, not of the profilers---which is itself a
characterization result a hardware study must respect (single-run accuracy
comparisons on such sequences are meaningless at sub-scatter resolution).

\section{Summary of the trust chain}

The system under measurement reproduces its vendor's published accuracy in
the characterized modes (141-run matrix); the instrumentation is
output-neutral (bit-identical where the workload is deterministic, and
within measured workload scatter elsewhere, across 192 variants); the noise
floor of all timing quantities is CoV 0.14\% at locked clocks
(\S\ref{sec:met-clocks}); and every classification threshold is
sensitivity-stressed (\S\ref{sec:met-sensitivity}). Every number in
Chapters~\ref{ch:characterization}--\ref{ch:synthesis} rests on this
chain.
